% ID: FIS-TER-GAS-006
% Titolo: Raddoppio di volume in un'espansione isoterma
% Disciplina: Fisica
% Area: Termologia
% Argomento: Gas perfetti
% Sottoargomento: Espansione isoterma reversibile
% Classe: Quarta liceo scientifico
% Difficolta: 4
% Tipo: problema_numerico
% Risultato: p_i = 9,98 * 10^4 Pa; p_f = 4,99 * 10^4 Pa; L = Q = 6,92 * 10^2 J
% Tempo_stimato: 12 min
% Competenze: equazione di stato, lavoro isotermo, primo principio, interpretazione del piano pressione-volume
% Prerequisiti: gas perfetti, trasformazione isoterma, logaritmi, primo principio della termodinamica
% Tag: termologia, gas_perfetti, isoterma, pressione, lavoro, calore, piano_pv
% Fonte: esempio_originale
% Autore: Riccardo Carli
% Licenza: CC-BY-SA-4.0
% Simulazione: ideal_gas_process

\begin{esercizio}
Un cilindro contiene \(0{,}40\,\mathrm{mol}\) di gas perfetto alla temperatura
costante di \(300\,\mathrm{K}\). Il gas compie un'espansione isoterma
reversibile passando dal volume
\[
V_i=1{,}00\cdot 10^{-2}\,\mathrm{m^3}
\]
al volume
\[
V_f=2{,}00\cdot 10^{-2}\,\mathrm{m^3}.
\]

Determina:
\begin{enumerate}
  \item la pressione iniziale e la pressione finale;
  \item il lavoro compiuto dal gas durante l'espansione;
  \item il calore assorbito dal gas.
\end{enumerate}
\end{esercizio}

\begin{soluzione}
Per un gas perfetto vale
\[
pV=nRT.
\]
La pressione iniziale e quindi
\[
p_i=\frac{nRT}{V_i}
=\frac{0{,}40\cdot 8{,}31\cdot 300}{1{,}00\cdot10^{-2}}
\simeq 9{,}98\cdot10^4\,\mathrm{Pa}.
\]

Poiche la trasformazione e isoterma, il prodotto \(pV\) resta costante. Il
volume finale e il doppio di quello iniziale, quindi la pressione finale e la
meta di quella iniziale:
\[
p_f=\frac{p_i}{2}
\simeq 4{,}99\cdot10^4\,\mathrm{Pa}.
\]

Il lavoro compiuto dal gas in un'espansione isoterma reversibile vale
\[
L=nRT\ln\frac{V_f}{V_i}.
\]
Pertanto
\[
L=0{,}40\cdot8{,}31\cdot300\ln 2
\simeq 6{,}92\cdot10^2\,\mathrm{J}.
\]

Per un gas perfetto l'energia interna dipende soltanto dalla temperatura. Essendo
la temperatura costante,
\[
\Delta U=0.
\]
Con la convenzione \(\Delta U=Q-L\), dal primo principio segue
\[
Q=L.
\]
Quindi il gas assorbe
\[
\boxed{Q\simeq 6{,}92\cdot10^2\,\mathrm{J}}.
\]
\end{soluzione}
